\section{神经网络为何需要非线性}

先做一个看似反常的实验。把许多线性层叠在一起：

\[
h_1=W_1x+b_1,\qquad
h_2=W_2h_1+b_2.
\]

展开便得

\[
h_2=(W_2W_1)x+(W_2b_1+b_2).
\]

无论继续叠多少层，最终仍只是一个 affine map。增加深度改变了参数化方式，却没有增加函数类别。神经网络获得复杂表示能力的关键，不是矩阵数量，而是在线性变换之间插入非线性。

\subsection{一层究竟做了什么}

前馈网络通常写成

\[
h_0=x,\qquad
a_\ell=W_\ell h_{\ell-1}+b_\ell,\qquad
h_\ell=\phi_\ell(a_\ell).
\]

affine 部分重新组合坐标、移动分界面；激活函数则弯折或截断空间，使下一层能对新的区域组合继续运算。以 ReLU 为例，

\[
\operatorname{ReLU}(z)=\max(0,z).
\]

单个 ReLU 形成一个 hinge。多个 hinge 在一维上可以拼出具有多个折点的连续分段线性函数；在高维中，不同激活模式把空间划分为许多区域，每个区域内部网络仍是 affine，但区域之间斜率不同。

例如函数

\[
f(x)=\operatorname{ReLU}(x)
-2\operatorname{ReLU}(x-1)
+\operatorname{ReLU}(x-2)
\]

在 \(x<0\) 时为 0，在 \(0<x<1\) 时上升，在 \(1<x<2\) 时下降，超过 2 后又回到 0。三个简单 hinge 已经构造出一个三角形局部特征。深层网络的优势在于复用中间特征，而不只是把更多 hinge 平铺在第一层。

\subsection{激活函数也决定梯度怎样流动}

sigmoid

\[
\sigma(z)=\frac1{1+e^{-z}}
\]

把实数映到 \((0,1)\)，但导数

\[
\sigma'(z)=\sigma(z)(1-\sigma(z))\le\frac14.
\]

上界 \(1/4\) 来自 \(s(1-s)\) 在 \(s=1/2\) 处取最大值。当 \(|z|\) 很大时，输出接近 0 或 1，导数也接近零。深层链式法则连续乘上这些小因子，早层梯度可能迅速消失。tanh 以零为中心，但同样会饱和。

ReLU 在正半轴导数为 1，缓解了正区域的饱和，却在负半轴导数为 0；若某单元长期落在负区，可能成为不再更新的 dying ReLU。GELU 等平滑激活改变了过渡方式，但不存在脱离初始化、归一化和任务语义的``永远最佳激活''。

输出层的映射尤其受任务约束：

\begin{itemize}
\tightlist
\item
  无界连续回归通常使用线性输出；
\item
  二元概率可使用一个 logit 配合稳定的 binary cross-entropy；
\item
  互斥多分类使用 logits 配合 softmax cross-entropy；
\item
  多标签任务使用多个独立 logits，而不是让标签在一个 softmax 中竞争。
\end{itemize}

``隐藏层激活''与``概率输出 link''承担不同职责，不能因为都叫 activation 就混用。

\subsection{表示能力不是训练与泛化保证}

通用逼近定理大意是：在适当条件下，足够宽、带合适非线性的一层网络可以任意逼近紧集上的连续函数。它说明某类函数在表示空间中存在，却没有回答：

\begin{itemize}
\tightlist
\item
  需要多少隐藏单元；
\item
  梯度方法能否找到该参数；
\item
  有限样本能否识别目标函数；
\item
  近似在训练分布之外是否可靠；
\item
  计算精度是否足以实现理论参数。
\end{itemize}

网络参数还存在大量对称性。交换同一隐藏层的两个单元并同步交换下一层对应权重，函数不变；对 ReLU 单元把入边乘 \(c>0\)、出边除以 \(c\)，函数也不变。两个参数向量相距很远，可能表示完全相同的函数。因此逐坐标比较权重通常不如比较输出、子空间或功能行为有意义。
